\subsection{A positive joint action on the actual rank-three tensor lattice}
The positive right matrices are already visible in \fid{F21}. The
following calculation strengthens that conclusion to the actual
ordinary left operators, tensor-product bar and full regular
commutant. Revised IV already gives a rank-three based construction
in Example~19.2, pp.~132--133; we do not claim priority or identify
our adapted basis with that particular designated basis.

\begin{finding}[\fid{F43}: actual positive actions, bar and graded paths]\label{f:positivejoint}
Put
\[
 \Lambda=\Z[v,v^{-1}],\quad \delta=v+v^{-1},\quad
 D=\delta^2,\quad h=D-1,
 \qquad S=\Lambda[(2h)^{-1}].
\]
The actual tensor lattice $M_\Lambda=(\Lambda^2\otimes\Lambda^2)^{\otimes3}$
has a free full-rank sublattice with basis $\mathcal B$ whose
change-of-basis determinant is $\pm2^{12}h^{72}$.
Thus $\mathcal B$ is a basis of the actual $M_S$ and is regular
at $v=1$ in characteristic different from $2,3$.
On this same basis, the actual right $P_1,P_2$, the ordinary left
$E_V,F_V,E_W,F_W,K_V^{\pm1},K_W^{\pm1}$, and the elementary
matrix units of the full commutant have nonnegative Laurent matrices.
Every vector of $\mathcal B$ is fixed by IV's actual ordinary
tensor-product bar. The commutant is
\[
 \mathcal C_S=\End_{A_S}(M_S)
 \simeq M_{20}(S)\times M_4(S)\times M_{16}(S)\times M_4(S),
\]
and its displayed matrix-unit basis acts by $0/1$ matrices.

For a nonempty right-generator word $w$ of length $L$, let $2s$
be its number of cyclic changes of letter. Then
\begin{equation}\label{eq:allwordpositive}
 \Tr_M(P_w)=20D^L+16D^{L-s}
              +4D^{L-2s}(D-2)^{2s}.
\end{equation}
Every entry of every operator sequence in the listed positive
operators counts paths in one explicit 64-state graded graph.
For explicitly written words this gives polynomial-time coefficient
evaluation and $O(L)$-bit positive witnesses, at fixed tower rank three.
\end{finding}

\paragraph{Left and right conventions.}
Right-generator matrices multiply in the written algebra order.
We also display a left operator by rows recording its coefficients
on basis vectors. For left algebra products these row matrices
reverse composition order: $\rho_{\rm row}(ab)=
\rho_{\rm row}(b)\rho_{\rm row}(a)$.
The path statements concern an explicitly ordered sequence of
applied operators. Reversing the written left-algebra word leaves
all the asserted positivity and complexity properties unchanged.

\begin{proof}
\par\smallskip\noindent\textbf{The ordinary tensor factors.}\quad
Use $\mathbf v_{ijk}$ for tensor monomials. The F21 matrix
\[
 B_3=\begin{pmatrix}
 v^2&v&1\\-1-v^{-2}&v&1\\-v^{-1}&-v^{-2}&\delta
 \end{pmatrix},\qquad \det B_3=h^2,
\]
acts on $(112,121,211)$ and on $(122,212,221)$, with its last
two rows interchanged in the second weight space. Write
\[
 b_0=\mathbf v_{111},\quad
 b_1=v^2\mathbf v_{112}+v\mathbf v_{121}+\mathbf v_{211},
\]
\[
 b_2=v^2\mathbf v_{122}+v\mathbf v_{212}+\mathbf v_{221},
 \qquad b_3=\mathbf v_{222}.
\]
Let $u_{+,1},u_{+,2}$ be the last two rows of $B_3$ in the first
weight space, and $u_{-,1},u_{-,2}$ the corresponding reversed
rows in the second. Set $u_{\pm,j}=g_\pm\otimes e_j$.
Let $W_3$ have basis $b_0,b_1,b_2,b_3$, $W_1$ basis $g_+,g_-$,
and $U$ basis $e_1,e_2$.
Here the subscripts are $\mathfrak{sl}_2$ highest weights;
the polynomial $\GL_2$ labels are $(3)$ and $(2,1)$, so
$W_1$ retains its determinant twist and polynomial degree three.
The actual ordinary Hecke action from \fid{F21} is
\[
 (S^2)^{\otimes3}=(W_3\otimes\epsilon)\oplus(W_1\otimes U),
 \qquad c_i|_\epsilon=\delta,
\]
\[
 c_1|_U=\begin{pmatrix}\delta&0\\1&0\end{pmatrix},\qquad
 c_2|_U=\begin{pmatrix}0&1\\0&\delta\end{pmatrix}.
\]
The determinant of the rank-eight basis change is $\pm h^4$.

The actual ordinary coproduct is
\cite[\S4.1, equation~(29), p.~22]{IV}
\[
 \Delta(E)=E\otimes K^{-1}+1\otimes E,
 \qquad\Delta(F)=F\otimes1+K\otimes F.
\]
Thus on a tensor word, $E$ changes a $2$ to $1$ with coefficient
$v^{-\sum\operatorname{wt}(i)}$, summed over later positions;
$F$ changes a $1$ to $2$ with coefficient
$v^{\sum\operatorname{wt}(i)}$, summed over earlier positions,
where $\operatorname{wt}(1)=1$, $\operatorname{wt}(2)=-1$.
Substituting the eight displayed basis vectors gives
\[
 Fb_0=b_1,\quad Fb_1=\delta b_2,\quad Fb_2=hb_3,\quad Fb_3=0,
\]
\[
 Eb_0=0,\quad Eb_1=hb_0,\quad Eb_2=\delta b_1,\quad Eb_3=b_2,
\]
\[
 Fg_+=g_-,\quad Fg_-=0,\quad Eg_-=g_+,\quad Eg_+=0.
\]
The $K$ eigenvalues are $v^{3-2j}$ on $b_j$ and $v^{\pm1}$
on $g_\pm$. All coefficients are nonnegative Laurent polynomials.

The bar assertion is also explicit. In the first weight space,
IV Figure~1, p.~9, gives the upper canonical vectors
\[
 c_{112}=\mathbf v_{112},\quad
 c_{121}=\mathbf v_{121}-v^{-1}\mathbf v_{112},\quad
 c_{211}=\mathbf v_{211}-v^{-1}\mathbf v_{121}.
\]
They are fixed by the actual tensor bar
\cite[\S6.2, Theorem~6.4, p.~34]{IV}.
The coordinates of $(b_1,u_{+,1},u_{+,2})$ in this basis are
\[
 \begin{pmatrix}h&\delta&1\\0&\delta&1\\0&1&\delta\end{pmatrix}.
\]
Every coefficient is bar invariant. The second weight space has
the same calculation with its last two rows reversed, and the two
endpoint vectors are canonical. Hence all eight adapted vectors
are fixed by the actual bar.

\par\smallskip\noindent\textbf{The full bimodule.}\quad
Regroup the rank-64 tensor module into the two ordinary rank-eight
factors and apply the preceding basis in each. On every Specht
factor $U\otimes U$ apply the F21 matrix
\[
 B_4=\begin{pmatrix}
 \delta&-1&-1&\delta\\0&-1&1&0\\
 2-D&-\delta&-\delta&D\\-D&\delta&\delta&D-2
 \end{pmatrix},\qquad \det B_4=8h^2.
\]
It gives $\epsilon_+\oplus\epsilon_-\oplus N(D-2)$,
while the mixed factors give $N(\delta)$, with
\[
 N(a):\quad
 P_1=\begin{pmatrix}D&0\\a&0\end{pmatrix},\qquad
 P_2=\begin{pmatrix}0&a\\0&D\end{pmatrix}.
\]
Consequently the actual bimodule is
\begin{align*}
 M_S={}&(W_3\otimes W_3\oplus W_1\otimes W_1)\otimes\epsilon_+
       \oplus(W_1\otimes W_1)\otimes\epsilon_-\\
 &\oplus(W_3\otimes W_1\oplus W_1\otimes W_3)\otimes N(\delta)
       \oplus(W_1\otimes W_1)\otimes N(D-2).
\end{align*}
The multiplicity ranks are $20,4,16,4$.
The $B_4$ change acts only on Specht factors, so the actual
ordinary left matrices remain the positive matrices just computed.
The right coefficients are $D=v^2+2+v^{-2}$,
$\delta=v+v^{-1}$ and $D-2=v^2+v^{-2}$.
The tensor bar for the direct-product quantum group is the product
of the two ordinary tensor bars after regrouping: its quasi-$R$
factors $\Theta_V,\Theta_W$ act on separate factors.
All $B_4$ coefficients are bar invariant, proving the full bar claim.

The rank-eight tensor-square basis change has determinant $h^{64}$,
by $\det(T\otimes T)=\det(T)^{16}$.
The four $B_4$ blocks contribute $(8h^2)^4=2^{12}h^8$.
All rows have coefficients in $\Lambda$, so their free span is a
sublattice with determinant $\pm2^{12}h^{72}$, and it becomes the
whole actual lattice over the stated $S$.

For the commutant, different generic simple sectors have zero
cross homomorphisms over the domain $S$.
This can also be seen from $(P_1-P_2)^2$, whose scalar values
on the scalar, $N(\delta)$ and $N(D-2)$ sectors are respectively
$0,D(D-1),4(D-1)$; $P_1$ separates the two scalar sectors.
The endomorphisms of each $N(a)$ are precisely scalars:
commuting a matrix $\left(\begin{smallmatrix}x&y\\z&t\end{smallmatrix}\right)$
with $P_1,P_2$ forces $y=z=0$ and $a(x-t)=0$, hence $x=t$.
The multiplicities give the stated product of matrix rings.
Each elementary left matrix unit changes only a multiplicity basis
vector, retaining its right factor, and therefore has a $0/1$
action on $\mathcal B$. Its usual multiplication constants are
also $0/1$. This is the actual regular commutant; identification
with every integral form of the dual FRT algebra is not asserted.

\par\smallskip\noindent\textbf{All right words.}\quad
In $N(a)$ put $u_1=(D,a)^T$, $u_2=(a,D)^T$, so $P_i=u_i e_i^T$.
If $w=i_1\cdots i_L$ has $t$ ordinary adjacent changes, multiplying
rank-one matrices gives
\[
 P_w=D^{L-1-t}a^t u_{i_1}e_{i_L}^T.
\]
Taking the trace introduces the final cyclic comparison.
The cyclic change count is even, say $2s$, and therefore
$\Tr_{N(a)}P_w=D^{L-2s}a^{2s}$.
Add the scalar-sector contribution $20D^L$, use
$\delta^{2s}=D^s$, and obtain \eqref{eq:allwordpositive}.
The empty word has trace $64$.

\par\smallskip\noindent\textbf{Positive paths and their size.}\quad
Replace an entry $\sum_d n_dv^d$ of a selected operator by $n_d$
directed edges of degree $d$ between its two basis states.
For example $P_1$ resets a rank-two block to state 1, with loop
degrees $-2,0,0,2$ and other-edge degrees $-1,1$ or $-2,2$;
$P_2$ resets it to state 2. Matrix multiplication proves, by
induction, that an operator sequence counts its color-prescribed
paths, with exponents equal to total degrees. Closed paths give
the trace. These are the actual conjugated operators, rather than
a pair chosen only to reproduce their traces.

For an explicitly written length-$L$ sequence, right generators
have at most four outgoing labeled edges, $E/F$ at most three,
and $K^{\pm1}$ or elementary commutant units at most one.
Every step has degree between $-3$ and $3$.
There are at most $4^L$ paths from a fixed start and at most
$64\cdot4^L$ closed paths, so witnesses and counts have $O(L)$
bits. Dynamic programming on 64 states and degrees in $[-3L,3L]$
uses polynomial bit time. This fixes tower rank three and uses
explicit words, not compressed binary powers or general
Kronecker/plethysm inputs.
\end{proof}

\paragraph{An exact graded model, and its summand boundary.}
For each matrix entry take a finite graded vector space with the
indicated homogeneous basis. These kernels define exact functors
on 64-tuples of finite graded vector spaces by direct sums and
tensor products. Composition is path concatenation. Their
Grothendieck-group action is exactly the displayed action on
$\Lambda^{64}$, and a word has the finite graded trace object
given by the sum of its diagonal kernel spaces.
The quadratic relation lifts by the tensor symmetry after a reset.
No coherent lift of all linear nonstandard relations or a
designated GCT categorification is claimed from the $K_0$ equality.

Closure under arbitrary direct summands has a proved limitation.
Keep just one copy of each right simple type, giving a faithful
rank-six module over the parameter ring (its special fiber need not
be faithful). At $v=1$, the two matrices $N(\delta)$
and $N(D-2)$ become identical. Thus every element of the actual
regular algebra acts identically on those two specialized blocks.
But the graded $P_1$ kernel has a direct summand consisting of
one degree-zero loop at state 1 of the $N(\delta)$ block, zero
elsewhere. Its action is a matrix unit supported on only that
branch and fails the preceding specialization condition.
It is outside the actual regular algebra. Hence closing this
semisimple matrix-functor model under all summands adds operations
beyond the nonstandard order, even with multiplicity copies removed.

The formal grading is also distinct from an arithmetic realization.
Closed-word traces in \eqref{eq:allwordpositive} have only even
$v$ powers, whereas individual $N(\delta)$ entries have odd powers.
Thus $v=q^{n/2}$ needs extra weight-one data at the entry level;
it is not automatically a geometric Tate construction.
Choosing $v=q^n$ uses a different, full-Tate normalization.

\paragraph{Two limits of simultaneous positivity.}
The actual factor swap $\tau$ fixes the first, third and fourth
rows of $B_4$ and negates its second row.
On the $\epsilon_-$ multiplicity space it is consequently
$-\mathrm{flip}$ on $W_1\otimes W_1$, of trace $-2$.
No basis of that summand makes $\tau$ nonnegative.
More generally, if $e_-$ is its actual projector, then
$\Tr(e_-\tau)=-2$, so $e_-$ and $\tau$ cannot both have
nonnegative matrices in any basis. The positive commutant cone
contains $e_-$. This agrees with IV's signed extension and is not
an objection to a based-module definition allowing that sign.

For $Q_i=DI-P_i$, negating the second vector of every rank-two
right block makes both $Q_i$ positive, preserves the actual bar
and leaves the ordinary left matrices positive.
But simultaneous positivity of both pairs $P_i,Q_i$ is impossible:
their off-diagonal entries are negatives of each other, so all
would have to vanish. This would make $P_1,P_2$ commute, whereas
their rank-two commutator has nonzero entry $Da$.

\paragraph{A strict positive center basis is also excluded.}
Over $R=\Q[v,v^{-1}]_{(v-1)}$, the actual rank-two center cannot
have an $R$-basis $\{1,b\}$ with
\[
 b^2=a(v)+c(v)b,\qquad a,c\in\N[v,v^{-1}].
\]
Its dual-number fiber forces the discriminant
$c(1)^2+4a(1)$ to vanish. Nonnegative coefficients imply
$a(1)=c(1)=0$, hence $a=c=0$ identically.
The generic algebra would then be $K[b]/b^2$, contradicting its
actual split reduced generic fiber. This excludes exactly the
stated unit-containing rank-two basis. F25's factor
$(v-v^{-1})^2(D-1)$ has negative expanded coefficients but is of
the factored kind allowed in VIII; the argument does not refute
that weaker positivity convention.

\paragraph{The prescribed-lattice gate remains.}
The adapted determinant is not a unit in $\Lambda$, so this basis
does not identify the original global tensor-monomial lattice.
It also does not lie in the ordinary upper lattice at infinity:
$b_1\otimes b_0$ has a $v^2$ tensor-monomial coefficient.
IV's paragraph after Conjecture~19.1 explicitly distinguishes its
desired nonstandard integral form and local lattice from these
ordinary ones. We have not identified our basis with those
prescribed nonstandard data or with Example~19.2.
Indeed, scaling an entire isotypic sector by the bar-invariant
unit $h^r$ preserves all the proved positive matrices and bar
properties over $S$, while changing the global and local lattices.
These properties therefore do not select VIII's prescribed
crystal reduction, duality and minimum-degree globalization.

% Separate paragraph for Part III.
\paragraph{\fid{F43} in Grothendieck's language.}
Over $\Spec S$ this is an explicit rank-64 vector bundle with
actual commuting algebra actions and a verified change of its
integral lattice. The full commutant is a product of matrix
algebras acting on its actual multiplicity bundles.
Separately, the finite graded kernels can be described by
$\mathbb G_m$-graded vector spaces on 64 points; their
Grothendieck group records the Laurent matrices and their
composition records paths. Closing that semisimple kernel
construction under every summand loses the coherent branch-gluing
condition, as the displayed specialization test proves.
The coordinate and $K_0$ descriptions make the positive action
concrete without identifying a formal quantum degree with a Weil
weight or specifying the dual FRT multiplication across degrees.
The actual bar and ordinary left action are now supplied in this
example; the designated nonstandard lattice, globalization and
geometric comparison remain separate requirements.

The equivariant-module and grading framework is SGA~3, I,
4.7.2--3, with cote~90, scan~17; the virtual $K_0$ and $\lambda$
framework is SGA~6, V, \S\S2--3 and cote~133, scan~18.
The displayed lattice, commutant and summand computations provide
the proof of this new application, rather than an archival attribution.
